\section{Rolle von Personas und UX-Strategien in der Praxis}
\label{sec:05-03_role-of-ux-methods}

Die Rolle von Personas sowie der wahrgenommene Einsatz von UX-Strategien im neuen Wiki wurden anhand geschlossener Skalenfragen, halb-offener Begründungen und offener Rückmeldungen erhoben. Ziel war es zu erfassen, inwieweit sich die Teilnehmenden in ihrer jeweiligen Rolle angesprochen fühlen und ob gestalterische Entscheidungen als nutzerzentriert wahrgenommen werden.

Die geschlossene Frage zur Passung des Wikis zur eigenen Rolle (GF-3) wurde überwiegend positiv bewertet. Der arithmetische Mittelwert lag bei 4{,}3 auf einer fünfstufigen Skala, der Median bei 4, bei einer Standardabweichung von 0{,}46. Drei Teilnehmende vergaben die höchste Skalenstufe, sieben Personen bewerteten die Rollenpassung mit einem Wert von 4; Bewertungen unterhalb dieser Stufe traten nicht auf (vgl.~Anhang). Damit ergibt sich eine hohe Zustimmung bei gleichzeitig geringer Streuung, was auf eine weitgehend einheitliche Wahrnehmung der Rollenadäquatheit hinweist.

Diese Einschätzung wird durch die qualitativen Begründungen (HF-3) weiter konkretisiert. Neun von zehn Teilnehmenden gaben an, dass alle für ihre Rolle relevanten Informationen im Wiki grundsätzlich vorhanden seien. Die Aussagen beziehen sich dabei sowohl auf operative Inhalte als auch auf übergeordnete Informationen zu Rollen, Verantwortlichkeiten und organisatorischen Zusammenhängen. Besonders häufig wurde hervorgehoben, dass das Wiki einen guten Überblick über Zuständigkeiten und Prozesse ermögliche, unabhängig davon, ob die jeweilige Rolle stark spezialisiert oder eher breit angelegt sei (vgl.~Anhang).

Sechs Teilnehmende betonten explizit, dass strukturierte Einstiegs- oder Übersichtsseiten die Rollenpassung unterstützen. Diese Seiten wurden als hilfreicher Ausgangspunkt beschrieben, um sich schnell zu orientieren oder neue Themen zu erschließen; insbesondere für weniger routinierte oder externe Nutzende stellten diese Einstiege eine zentrale Unterstützung dar. Gleichzeitig äußerten drei Teilnehmende den Wunsch nach einer noch stärkeren Rollendifferenzierung, etwa durch klarere Kennzeichnungen, welche Inhalte primär für bestimmte Rollen oder Erfahrungsstufen relevant sind. Diese Rückmeldungen beziehen sich nicht auf fehlende Inhalte, sondern auf eine weitergehende Feinjustierung der bestehenden Struktur (vgl.~Anhang).

Auffällig ist, dass keine Person angab, sich durch das Wiki in der eigenen Rolle nicht angesprochen oder falsch adressiert zu fühlen. Die bewusst generisch gehaltene Struktur wurde von allen Befragten akzeptiert. Die qualitative Analyse zeigt, dass diese Offenheit als Vorteil wahrgenommen wird, solange zentrale Rolleninformationen leicht auffindbar bleiben. Die Rollenpassung wird damit weniger über individualisierte Oberflächen als über eine klare Informationsarchitektur und transparente Zuständigkeiten hergestellt.

Der wahrgenommene Einsatz von UX-Strategien wurde sowohl implizit über die Interviews als auch explizit über eine geschlossene Skalenfrage zur Nutzerorientierung (GF-4) erfasst. Diese Frage wurde mit einem arithmetischen Mittelwert von 4{,}2 und einem Median von 4 beantwortet, bei einer Standardabweichung von 0{,}6. Acht Teilnehmende stimmten der Aussage zu, dass das Wiki erkennbar nutzerzentriert gestaltet sei, zwei Personen wählten eine neutrale bis leicht zustimmende Antwortoption; ablehnende Bewertungen traten nicht auf (vgl.~Anhang).

In den offenen Rückmeldungen (OF-2 und OF-3) äußerten sieben von zehn Teilnehmenden explizit die Vermutung, dass bei der Entwicklung des neuen Wikis UX-Methoden oder zumindest eine nutzerzentrierte Herangehensweise eingesetzt worden seien. Als Indikatoren hierfür nannten sie unter anderem die klare Struktur, die reduzierte visuelle Gestaltung, konsistente Benennungen sowie die Orientierung an konkreten Nutzungsszenarien. Drei Teilnehmende gaben an, den Einsatz formaler UX-Methoden nicht eindeutig beurteilen zu können, beschrieben die Gestaltung jedoch dennoch als \glqq durchdacht\grqq{} oder \glqq praxisnah\grqq{} (vgl.~Anhang).

Zusammenfassend zeigen die Ergebnisse, dass das neue Wiki von der Mehrheit der Teilnehmenden als gut zur eigenen Rolle passend wahrgenommen wird. Die positiven Bewertungen beziehen sich vor allem auf die Vollständigkeit der Inhalte, die Unterstützung durch strukturierte Einstiegsseiten sowie die generische, aber nachvollziehbare Informationsarchitektur. Gleichzeitig weisen einzelne Rückmeldungen auf Potenziale zur stärkeren Sichtbarmachung rollenrelevanter Inhalte hin. Der Einsatz von UX-Strategien wird überwiegend als erkennbar wahrgenommen und mit konkreten Gestaltungsmerkmalen begründet. Diese Ergebnisse bilden die empirische Grundlage für die Identifikation übergreifender Erfolgsfaktoren und Muster im folgenden Kapitel.
